\documentclass[runningheads]{llncs}
\usepackage[T1]{fontenc}
\usepackage{graphicx}
\usepackage{booktabs}
\usepackage{amsmath}
\usepackage{subcaption}
\usepackage{hyperref}

\begin{document}

\title{Sneakie: Learning to Play the Snake Game in All Its Forms}
\titlerunning{Sneakie: Deep RL on Snake}

\author{JAMAL Adonis \and  KAPOTOS Fotios \and MARTINI Jean-Vincent }
\institute{CentraleSupélec}
\maketitle

% ── Abstract ──────────────────────────────────────────────────────────────────
\begin{abstract}
We present a comparative study of deep reinforcement learning methods on Snake,
a grid-world game that exhibits many of the challenges found in real sequential
decision-making problems: sparse rewards, partial observability, long-horizon
planning, and increasingly complex environment dynamics.
We develop a Snake environment that supports multiple food types (gold, silver, and poison), static and dynamic obstacles, and three distinct observation spaces: a 15-dimensional hand-crafted feature
vector, an egocentric convolutional window, and a full-grid CNN observation.
Using DQN and its extensions (Double DQN, Dueling networks, frame stacking), we study the ability of our agents to solve, generalize and adapt to a range of different environments, from easy to complex.
Our results show that CNN-based agents consistently outperform hand-crafted
MLP agents in complex environments.
You can find our extensive results here: \url{https://wandb.ai/fotiskapotos/rl-snake?nw=nwuserfotiskapotos} and our code here: \url{https://github.com/JeanJNMV/MDS-RL_SnakeAgent}.
\end{abstract}

% ── 1. Introduction ───────────────────────────────────────────────────────────
\section{Introduction}
\label{sec:intro}

The game of Snake is a simple environment that captures key challenges in
reinforcement learning: the agent must collect food to grow while avoiding
collisions with walls, obstacles, and its own body.
To increase the richness of the environment, we introduce a \emph{hard mode}
with multiple food types (gold, silver, poison) and dynamic obstacles.

We study how different state representations and training strategies impact
performance across environments of increasing complexity.
Notably, a standard DQN agent~\cite{mnih2015human} quickly reaches very good performance on the classic Snake game \cite{wei2018autonomous}, effectively hitting a performance ceiling; we therefore designed a substantially harder variant of the environment and shifted our focus to understanding how different architectures and training techniques cope with this increased complexity.

\noindent\textbf{Main contributions.}
\begin{itemize}
  \item A configurable Snake environment with multiple observation spaces,
        reward functions, and a hard mode with dynamic obstacles.
  \item A comparison of hand-crafted features and CNN-based representations.
  \item An evaluation of DQN extensions (Double DQN, Dueling networks,
        frame stacking) in complex environments.
\end{itemize}
 
\section{Environment Design}
\label{sec:env}

\paragraph{State and observations.}
The environment is a discrete $H \times W$ grid (typically $16 \times 16$),
where each cell corresponds to empty space, snake body/head, gold, silver,
poison, or obstacle. We consider three observation modalities:\\
(i) \textbf{MLP features}: a 15-dimensional binary vector encoding collision
risks (straight/left/right), current direction, and relative positions of
nearest food and poison;\\
(ii) \textbf{Egocentric CNN}: a $C \times k \times k$ tensor centered on the
head with $C=6$ channels (body, head, gold, silver, poison, obstacles),
capturing local spatial structure;In our case we choose $k=6$.\\
(iii) \textbf{Full-grid CNN}: a $C \times H \times W$ tensor providing global
context at higher dimensional cost.

\paragraph{Dynamics and actions.}
At each step, the agent selects one of four absolute actions (up, right, down,
left). Reverse moves can be disallowed, in which case the snake continues
forward. The dynamics follow standard Snake rules: the head moves, the tail is
removed if no food is consumed, the snake grows when eating gold or silver,
and shrinks when consuming poison.

\paragraph{Rewards.}
Rewards are additive: gold $+r_g$, silver $+r_s$, poison $+r_p \leq 0$, death
$r_d < 0$, and a step reward $r_{\text{step}}$. When no food is consumed, we
add potential-based shaping \cite{ng1999policy}:
\begin{align}
r_t \mathrel{+}= \lambda (d_{t-1} - d_t) + \mu (b_t - b_{t-1}),
\end{align}
where $d_t$ is the Manhattan distance to the nearest non-poison food and $b_t$
the distance to the nearest body segment. The first term promotes food-seeking,
while the second discourages self-collisions.

\paragraph{Obstacles and settings.}
We consider static obstacles (fixed), random obstacles (sampled at reset), and
dynamic obstacles (3-cell walls oscillating along a fixed axis). Food is never
spawned in occupied or swept regions. We define two regimes:
\textbf{easy} ($16 \times 16$, one gold item, no obstacles, static dynamics)
and \textbf{hard} ($16 \times 16$, multiple foods including poison, static,
random, and dynamic obstacles).
% ── 3. Methods ────────────────────────────────────────────────────────────────
\section{Methods}
\label{sec:methods}

All agents share the same DQN~\cite{mnih2015human} training loop: $\epsilon$-greedy exploration with
linear annealing, experience replay (buffer size 100\,000), and a target
network updated every $C$ steps.

\paragraph{Architectures.}
We consider three state representations.
(i) \textbf{MLP (baseline):} a 3-layer fully-connected network
($256$--$128$--$3$) operating on a 15-dimensional hand-crafted feature vector,
which directly encodes collision risks, direction, and relative food/poison
positions.
(ii) \textbf{Egocentric CNN:} two convolutional layers followed by two
fully-connected layers processing a $6 \times k \times k$ window centered on
the head, enforcing locality and translation equivariance.
(iii) \textbf{Full-grid CNN:} the same architecture applied to the full
$C \times H \times W$ board, providing global spatial context at higher
dimensional cost.

\paragraph{DQN extensions.}
We use \textbf{Double DQN}~\cite{van2016deep} to reduce overestimation bias by decoupling action
selection and evaluation:
\begin{equation}
  y_t = r_t + \gamma\, Q\bigl(s_{t+1},\,\arg\max_a Q(s_{t+1}, a;\,\theta);\,\theta^-\bigr).
\end{equation}
We also use \textbf{Dueling networks}~\cite{wang2016dueling}, decomposing $Q(s,a)$ into a state-value
$V(s)$ and an advantage $A(s,a)$ as
$Q(s,a) = V(s) + \bigl(A(s,a) - \overline{A}(s)\bigr)$, which improves learning
in states where actions have similar value.

\paragraph{Temporal encoding.}
For CNN agents, we stack the last $n=4$ observations into a
$(4C) \times H \times W$ tensor. This provides implicit velocity information
for dynamic obstacles and mitigates partial observability. We find that $4$ frames is the best compromise
between computational effeciency and good state representation. 

% ── 4. Experiments ────────────────────────────────────────────────────────────
\section{Experiments and Results}
\label{sec:exp}

\subsection{MLP agent on the easy environment}
\label{sec:exp_baseline}
We start by training an MLP agent on the easy environment.
The agent quickly learns to collect food and avoid collisions.
We then study the impact of reward design by introducing a step penalty
and distance-based shaping to accelerate learning.
While these modifications slightly improve convergence speed, the overall
performance gains remain marginal.
We also conduct a brief analysis of reward scaling by multiplying all rewards
by a constant factor. This has limited impact on the final performance,
but can affect training stability and the speed of convergence.

\begin{figure}[h]
  \centering
  \begin{subfigure}{0.35\linewidth}
    \centering
    \includegraphics[width=\linewidth]{../poster/img/MLP rewards.png}
    \caption{Baseline vs.\ step penalty vs.\ distance shaping}
  \end{subfigure}
  \hfill
  \begin{subfigure}{0.35\linewidth}
    \centering
    \includegraphics[width=\linewidth]{../poster/img/reward magnitude_scaling.png}
    \caption{Effect of reward magnitude scaling}
  \end{subfigure}
  \caption{MLP agent snake length over training (5\,000 episodes, $16\!\times\!16$ grid).}
\end{figure}

\subsection{CNN agents on the easy environment}
We next replace the MLP with convolutional agents operating directly on
grid observations, considering both egocentric-window and full-grid inputs
(Fig.~\ref{fig:cnn_comparison}).

The egocentric CNN learns faster in the early stages of training
(Fig.~\ref{fig:ego_vs_full}), benefiting from its focused local representation.
However, it converges to a slightly lower final performance than the full-grid
CNN, which has access to global spatial information.

We also compare CNN and MLP agents (Fig.~\ref{fig:mlp_vs_cnn}). While the CNN
converges more slowly, it achieves slightly better final performance.

\begin{figure}[h]
  \centering
  \begin{subfigure}{0.44\linewidth}
    \centering
    \includegraphics[width=\linewidth]{../poster/img/Ergovsfull.png}
    \caption{Egocentric vs.\ full-grid CNN}
    \label{fig:ego_vs_full}
  \end{subfigure}
  \hfill
  \begin{subfigure}{0.44\linewidth}
    \centering
    \includegraphics[width=\linewidth]{../poster/img/train_simple.png}
    \caption{MLP vs.\ CNN}
    \label{fig:mlp_vs_cnn}
  \end{subfigure}
  \caption{Performance comparison of architectures on the easy environment.}
  \label{fig:cnn_comparison}
\end{figure}

\subsection{Results on the Hard Environment}
\subsubsection{Adaptation of Pre-trained Agents}
We evaluate how agents trained in simpler settings transfer to the hard
environment. Performance drops significantly for the MLP, particularly in
the presence of dynamic obstacles, where both reward and survival degrade
and the loop rate increases sharply. This suggests that hand-crafted features
fail to capture the additional complexity and non-stationarity of the environment.

In contrast, CNN agents exhibit substantially better robustness. While their
performance also decreases under moving obstacles, the degradation remains
limited and loop behaviours are less frequent. This indicates that spatial
representations enable better adaptation to complex dynamics and obstacle
interactions.

Overall, the results highlight a clear gap in generalization: CNN-based agents
transfer more effectively to the hard regime, whereas MLP agents struggle to
cope with increased environmental complexity.

\begin{table}[h]
\centering
\small
\begin{tabular}{lcccc}
\toprule
\textbf{Agent} & \textbf{Mean length} & \textbf{Mean reward} & \textbf{Mean foods} & \textbf{Loop rate} \\
\midrule
MLP (static hard)   & 13.2 & 10.8 & 10.2 & 0.05 \\
MLP (moving obs.)   &  9.1 &  6.7 &  6.4 & 0.27 \\
CNN (static hard)   & 15.6 & 12.9 & 12.5 & 0.02 \\
CNN (moving obs.)   & 14.8 & 11.7 & 11.3 & 0.08 \\
\bottomrule
\end{tabular}
\caption{Performance of MLP and CNN agents in the hard environment.}
\label{tab:hard_results}
\end{table}
\subsubsection{Training CNN Agents on the Hard Environment}

We finally train a CNN agent directly on the hard environment.
We incorporate the full set of stabilisation techniques introduced earlier,
including Double DQN, Dueling DQN, and frame stacking, and tune the reward
parameters to improve convergence and stability.

\begin{figure}[h]
  \centering
  \includegraphics[width=0.5\linewidth]{../poster/img/finalfinal.png}
  \caption{Training performance of the CNN agent on the hard environment.}
\end{figure}


% ── 6. Conclusion ─────────────────────────────────────────────────────────────
\section{Conclusion}
\label{sec:conclusion}

We presented a comparative study of Deep Q-Network variants on a configurable Snake environment of increasing complexity. Our results show that while hand-crafted MLP agents perform well in simple settings, they fail to scale to more complex and dynamic environments. In contrast, CNN-based agents, especially when combined with Double DQN, Dueling architectures, and frame stacking, achieve stronger performance and exhibit improved generalization.

Overall, these results highlight the importance of spatial representations and stabilisation techniques for reinforcement learning problems involving rich dynamics and partial observability.

\paragraph{Limitations.}
This study has several limitations. First, experiments are conducted on relatively small grid sizes ($16 \times 16$), and scaling to larger environments remains unexplored. Second, generalization is only evaluated across variations of the same environment, and not across fundamentally different tasks. Third, performance is sensitive to hyperparameter choices and reward shaping, which may affect reproducibility and obscure intrinsic method comparisons. Fourth, due to computational constraints, results are based on single training runs rather than multiple seeds, limiting the statistical robustness of our conclusions. Finally, our analysis is purely empirical and does not provide theoretical guarantees on stability or convergence.

\bibliographystyle{splncs04}
\bibliography{bibliography}

\end{document}
